Carrying an open container that is filled to the brim is a task in which speed and safety pull in opposite directions: every acceleration of the container tilts and excites the free surface, and the surface has only a few millimetres of freeboard before liquid is lost. The practical objective, as stated in the brief for this study, is not merely to avoid spilling (that is easy if one moves slowly) but to be \emph{the fastest spill-free transport}.

A controller that plans fast motions needs a prediction of the surface. The dominant choice is an equivalent mechanical model, a pendulum or a mass-spring-damper whose parameters are fitted to the liquid \cite{capolupo2025equivalent,guagliumi2021simple,leva2022sloshing}. Such models underlie slosh-free trajectory optimisation \cite{muchacho2022solution}, feedforward rest-to-rest manoeuvres \cite{toth2024rapid}, time-optimal planning with a bound on the surface height \cite{ferrari2026time} and real-time MPC \cite{medico2026model,hynninen2026emergency}. A learned dynamics model is an alternative that has been used for spill-free transport \cite{abderezaei2024clutter}; it can absorb higher modes and nonlinear effects, and with a history encoder it can adapt to a liquid whose properties are not known in advance \cite{kumar2021rma,evans2022context}.

What is missing is a like-for-like comparison: the same plant, optimiser, spill criterion and tuning rule, with only the slosh model exchanged. This paper reports such a comparison at small scale. The question is: \emph{does a learned slosh predictor inside sampling-based MPC enable faster spill-free transport than MPC with a fitted pendulum or spring-mass model, across liquid parameters and external pushes?}

\textbf{What we did.} We built a planar cart-and-container simulation with a reduced-order two-mode nonlinear slosh ground truth, a time-minimising MPPI controller with a pluggable slosh model, three fitted parametric models (pendulum, spring-mass, two-mode spring-mass), a zero-vibration input-shaping baseline, and a learned predictor with and without latent context inference. Each system's single safety parameter is tuned by one validation rule. We test four hypotheses fixed before the experiments:
\begin{itemize}
\item \textbf{H1}: on the nominal liquid the learned-model MPC is faster, spill-free, than the fitted-pendulum MPC.
\item \textbf{H2}: on held-out liquids with pushes, the learned model with a history encoder spills less than the fixed-parameter pendulum MPC at equal transport time.
\item \textbf{H3}: the learned advantage vanishes when the true slosh is a single linear mode.
\item \textbf{H4}: the sensor set changes performance, and force/torque sensing is the most informative single addition to cart-only sensing.
\end{itemize}

\textbf{What we found.} H1 holds against the pendulum and the single-mode spring-mass model, but the learned controller is \emph{not} the fastest system on the nominal liquid: a fitted two-mode spring-mass model and open-loop input shaping are faster there. H2 holds only in the equal-time sense; at each system's own tuned margin the slower pendulum controller spills less. H3 holds. The first half of H4 holds and the second half is not supported. The study is a simulation with an invented reduced-order liquid, five seeds, and reimplemented baselines; Section~\ref{sec:limitations} states what this does and does not show.
